% RCV reorganization of the Method, around the reactive advantage as a value of information.
% Draft by worker-B; matches the notation of method.tex (x_t=(o_t,a_{t-m:t-1}), Q(x,c,k), kappa)
% and cites the labels in theory.tex (prop:bookkeeping, prop:optimism, lem:fullchunk, cor:curriculum,
% rem:confound, rem:sign). NOT a drop-in replacement -- a proposed reframing; integrate or discard.
\section{Method}
\label{sec:method}

Adaptive commitment asks, at each query, whether to keep executing a plan or to stop and replan.
We make that question a single scalar the critic can estimate and the deployment rule can read: the
\emph{reactive advantage}, the value of committing a prefix over replanning now. Its sign is a value
of information --- positive exactly where the current observation has lost what the plan still carries.
The method is the critic that estimates this quantity without the bias that closed-loop demonstrations
otherwise inject, an actor that improves the plan against it, and a deployment rule that reads it.

\subsection{Setup and the reactive advantage}
\label{sec:method:setup}

A pretrained policy $\pi_\theta$ maps an observation $o$ to a chunk of $H$ actions; $c=a_{1:H}$ is a
chunk and $c_{1:k}$ its first $k$ actions. Executing $c_{1:k}$ open loop and then querying again is a
temporally extended action of duration $k$, so values across commitments must be compared with
semi-Markov bookkeeping~\citep{sutton1999options,bradtke1994smdp},
\begin{equation}
\label{eq:rcv-smdp}
Q^{\pi,\kappa}(x, c, k) \;=\; \mathbb E\Big[\textstyle\sum_{j<k}\gamma^{j} r_j \;+\; \gamma^{k}\,
V^{\pi,\kappa}(x_k)\Big],
\end{equation}
where $x$ is the critic's input and $\kappa$ the commitment rule. Charging one discount per decision
rather than per elapsed step distorts the comparison by a $k$-dependent term whose sign follows the
reward convention, so a duration-blind value induces a preference over commitments that is an artifact
of the value scale, not of the task (Prop.~\ref{prop:bookkeeping}, Rem.~\ref{rem:sign}); the backup of
\eqref{eq:rcv-smdp} removes it.

Write $V_{\mathrm{react}}(x)=Q(x,c,1)$ for the value of replanning at the next step and
$Q_{\mathrm{commit}}(x,c,k)=Q(x,c,k)$ for committing $k$. The selection signal is their difference,
\begin{equation}
\label{eq:rcv-advantage}
A(x,c,k) \;=\; Q_{\mathrm{commit}}(x,c,k) \;-\; V_{\mathrm{react}}(x).
\end{equation}
Both terms are the same value functional evaluated on the same belief over latent state induced by
$x$; they differ only in the information the executed actions were chosen with --- the plan $c$ was
drawn when the branch was observable, whereas replanning decides from $x$ alone. Hence $A$ is the
value of the information the current input lacks: it is nonpositive when $x$ is a sufficient statistic
of the state, since replanning can then reproduce the plan and interrupting a commitment is weakly
optimal~\citep{sutton1999options}, and it is strictly positive only when the observation has been
garbled relative to when the plan was made, so committing carries information forward that replanning
would discard. The rest of the method is what it takes to estimate $A$ so that the sign of the
estimate matches the sign of the truth.

\subsection{An honest reactive advantage}
\label{sec:method:critic}

Three properties of closed-loop demonstration data corrupt a naive estimate of $A$, each in a definite
direction, and each is removed by a modeling choice rather than a loss term. We state each as the bias
it removes.

\paragraph{History conditioning is forced, not added.}
The input is not a modeling preference; it is what makes a positive $A$ representable at all. Suppose
the critic's input $x$ is a sufficient statistic of the state, as it is when $x=o_t$ and the
environment is fully observed. Then re-planning at any intermediate step has access to at least the
information the commitment was chosen with, so the interrupted option is weakly better than the
committed one~\citep{sutton1999options}: $V_{\mathrm{react}}(x)\ge Q_{\mathrm{commit}}(x,c,k)$ for every
$k$, that is $A\le 0$ everywhere, with equality only for a self-consistent optimal policy. An
observation-conditioned critic therefore cannot express a strict preference for a longer commitment; at
best it ties, and any strict preference it does report is estimation error. Every force that breaks the
tie inside this class --- policy error, bootstrap optimism --- points short. So a critic that is to
price commitment at all must leave the class in which its input determines the state.

It leaves that class exactly when the demonstrator conditioned on something the input does not carry.
A teleoperator holds a plan across a stretch of motion and acts from it while the scene alone no longer
reveals it; the chunk drawn at the informative moment carries that latent forward, whereas re-planning
re-conditions on the impoverished present and discards it. Conditioning on
\begin{equation}
\label{eq:rcv-input}
x_t=(o_t,\, a_{t-m:t-1})
\end{equation}
restores the latent through the actions already executed since the last query --- the past commitment
is itself evidence about what the demonstrator was doing --- and a delayed observation $o_{t-\Delta}$
serves the same role more cheaply. With it, $A>0$ is representable precisely where the plan carries
information the current input has lost, which is the content of the sign law in
Section~\ref{sec:method:setup}.

Two consequences worth stating. First, the conditioning belongs on the \emph{critic}: it is trained
offline and unconstrained by inference latency, while the policy stays short context. Conditioning the
policy on history instead lowers its training loss but worsens its rollout, because it overfits temporal
correlations that break under its own induced state distribution~\citep{park-bc-mystery}; the same
elimination result also says the alternative fix --- simply lengthening the policy's
context~\citep{zeng2026openloop} --- is unavailable to a deployed VLA. Second, the same result names the
failure mode we must avoid: history helps only insofar as $x$ carries the latent the observation lost,
rather than an action-history shortcut the critic can copy. That distinction is measurable, and the
reactive-map diagnostic of Section~\ref{sec:method:diagnostic} is the test.

\paragraph{Interventional composition removes the leakage in $Q_{\mathrm{commit}}$.}
The committed value must not be estimated by regressing outcomes on the executed chunk. An event
revealed inside the window causes both the recorded actions and the outcome, so conditioning the value
on the chunk shifts the latent posterior toward the states where that chunk succeeded --- the critic
is credited with the demonstrator's foreknowledge. Composing one-step backups is necessary but not
sufficient: bootstrapping each transition from its own recorded successor preserves the correlation.
We compose while resampling the successor among transitions at the same decision point, holding the
evaluated tail fixed,
\begin{equation}
\label{eq:rcv-interv}
y_k(x_t, c) \;=\; r_t \;+\; \gamma\, \mathbb E_{x' \sim \mathcal{D}(\cdot \mid x_t)}
\big[\, Q_{\bar\phi}(x', c_{2:k}, k-1) \,\big],
\qquad k > 1,
\end{equation}
with $\mathcal D(\cdot\mid x_t)$ the empirical successor distribution at decision points matching
$x_t$ and $\bar\phi$ the target parameters. The backup then estimates $Q(x,\mathrm{do}(c_{1:k}))$,
the quantity open-loop execution faces; it holds the posterior at the belief the input warrants
instead of the belief the chunk reveals. Under open-loop consistent data the two coincide~\citep{dqc};
the intervention is what buys that independence when the data is closed loop.

\paragraph{An honestly priced requery removes the optimism in $V_{\mathrm{react}}$.}
The leaf of the recursion is the value of querying the policy again, and it must be the value of the
policy that will resume control. Two leaves fail here, in the same direction. Bootstrapping from
demonstration returns credits the requery with the demonstrator's competence. Bootstrapping from a
maximum over sampled candidates --- the EMaQ-style operator that prior adaptive-chunking critics use,
taking a joint $\max$ over $N$ candidate chunks and $H$ prefixes~\citep{acsac} --- credits it instead
with the largest estimation error among $NH$ options, since $\mathbb E[\max] \ge \max \mathbb E$ and the
gap widens with the number of options. Because the true values of different commitments coincide in the
idealized limit, that maximum is taken over what is mostly estimation noise. Either optimism enters the
prefix comparison with weight $\gamma^{k}$ and therefore biases the choice toward short commitments
regardless of the reward convention (Prop.~\ref{prop:optimism}). We keep the deployment-consistency of
the second (the leaf is the policy that will actually resume) while removing the maximum, using the
deployed policy's own expectation,
\begin{equation}
\label{eq:rcv-leaf}
y_1(x_t, c) \;=\; r_t \;+\; \gamma\, \mathbb E_{x'}\,\mathbb E_{c'\sim\pi_\theta(\cdot\mid o')}
\big[\, Q_{\bar\phi}\big(x', c', \kappa(x', c')\big) \,\big],
\end{equation}
so $V_{\mathrm{react}}$ is a fixed point of the process that will run, not of an idealized one. The
chunks scored during training are the recorded ones; only the leaf involves the current policy, one
step deeper where the state is more informative.

The critic is a causal transformer over the chunk on frozen backbone features, emitting all $H$ prefix
values in one pass, trained by regression on $\{y_k\}_{k=1}^{H}$ from
\eqref{eq:rcv-interv}--\eqref{eq:rcv-leaf} with a target network. With these three choices the sign of
the estimated $A$ tracks the truth: nonpositive where the input is sufficient, positive where a plan
survives an occlusion the observation does not.

\subsection{Extraction over the joint action--horizon space}
\label{sec:method:actor}

The reactive advantage prices commitments for a \emph{fixed} chunk, but the chunk is not fixed: the
policy is what we improve. Two couplings make action and horizon inseparable. The optimal action
depends on how long it will be committed --- a reactive first action where a branch is imminent, a
coherent multi-step plan where none is --- and the horizon the selector chooses depends on the chunk,
since a chunk's prefix values determine $\kappa(x,c)$. Improving the two separately mis-specifies each.

\paragraph{Why the full-chunk objective is not enough.}
Improving the policy on the full-chunk value $Q(x,c,H)$ is stable and lower-bounds deployed value
(Lem.~\ref{lem:fullchunk}), but it optimizes a plan of length $H$ everywhere, including states where the
selector will commit only $k\ll H$. There the tail is discarded, and because a chunk model trades
capacity across its actions, optimizing that tail can compromise the very first action --- the one that
is actually executed and that should react. The full-chunk objective is the right target where the
optimal commitment is long and the wrong one where it is short.

\paragraph{The joint objective.}
The value deployment realizes for proposing $c$ at $x$ is the value at the length actually executed,
$Q(x, c, \kappa(x,c))$, not $Q(x,c,H)$. We therefore improve the policy against the deployed value
directly, reinforcing the prefix that will run, weighted by its advantage at that length,
\begin{equation}
\label{eq:rcv-joint}
\max_\theta\; \mathbb E_{(x, c)\sim\mathcal D}\Big[\,
w(x,c)\,\log \pi_\theta\big(c_{1:\kappa(x,c)} \mid o\big)\Big],
\qquad
w(x,c)=\exp\!\Big(\beta^{-1}\big[\,Q_\phi\big(x, c, \kappa(x,c)\big)-V_{\mathrm{react}}(x)\,\big]\Big),
\end{equation}
with $V_{\mathrm{react}}(x)=\mathbb E_{c'\sim\pi_\theta}[Q_\phi(x,c',\kappa(x,c'))]$ and $\kappa$ from
\eqref{eq:rcv-kappa}. This is optimization over the joint action--horizon space: a chunk is reinforced
on exactly the prefix deployment will execute, and by its value at that length. Where the state is
reactive, $\kappa=1$ and only the first action is improved, toward a good reaction; where it admits a
plan, $\kappa$ is long and the whole coherent chunk is improved. The full-chunk objective (advantage-weighted regression at $k=H$) is the special case $\kappa\equiv H$; \eqref{eq:rcv-joint} refines it per state.

\paragraph{Why it does not collapse to short.}
Reinforcing only short prefixes, independent of value, would leave the full-chunk value unchanged and
never lengthen commitments. The joint weight avoids this because it is a \emph{value} weight: where a
long plan is genuinely good, the chunks with high $Q(x,c,H)$ are the ones the selector commits at
$\kappa=H$, so \eqref{eq:rcv-joint} reinforces their full length and the policy produces more of them.
Commitments lengthen exactly where lengthening carries value, which is the curriculum improvement
predicts (Cor.~\ref{cor:curriculum}), while reactive states keep a first action optimized to react.

\paragraph{The coupling, and stability.}
Because $\kappa$ reads $Q_\phi$, which depends on $\pi_\theta$, \eqref{eq:rcv-joint} is nonstationary:
the horizon that defines the objective moves as the policy improves. We compute $\kappa$ from a target
critic $Q_{\bar\phi}$ and treat it as a stop-gradient selection, and we alternate a critic evaluation
step \eqref{eq:rcv-interv}--\eqref{eq:rcv-leaf} with the actor step \eqref{eq:rcv-joint}. The honesty of
the critic is a precondition rather than a detail: if $A$ is contaminated by the leakage of
Section~\ref{sec:method:critic}, $\kappa$ selects the wrong length and the joint objective reinforces
the wrong prefix, so the interventional and requery corrections are what make joint extraction safe.

\subsection{Deployment}
\label{sec:method:deploy}

At each query the policy emits one chunk, the critic scores its prefixes in the same forward pass, and
the system executes the longest admissible prefix within $\epsilon$ of the best,
\begin{equation}
\label{eq:rcv-kappa}
\kappa(x, c) \;=\; \max\Big\{\,k \in \mathcal K \;:\; A(x,c,k) \;\ge\; -\epsilon \,\Big\}
\;=\; \max\Big\{\,k \in \mathcal K \;:\; Q_\phi(x, c, k) \;\ge\;
\max_{k'} Q_\phi(x, c, k') - \epsilon \,\Big\}.
\end{equation}
The menu $\mathcal K\subseteq\{1,\dots,H\}$ is a design choice reported with every result; a coarse
menu costs little, since neighbouring commitments have close values, and keeps results comparable to
prior work that adapts over a small set. Breaking ties toward long commitments, not short, is what
yields the curriculum: constant requerying re-injects policy error and destroys the plan continuity
the demonstration relied on. Apart from the menu, $\epsilon$ is the only free parameter, and no cost
term is added to the reward --- the pressure to commit is the value of information in $A$, not a price.

\subsection{What the critic measures}
\label{sec:method:diagnostic}

The reactive advantage is a per-state quantity, and its magnitude is the non-Markovian content of the
state: the mutual information $I(s;c\mid o)$ between the demonstrator's chunk and the latent given the
observation, cashed out in value units. It equals the gap between the leaky and interventional targets,
$Q_{\mathrm{reg}}-Q_{\mathrm{syn}}$, which we report per frame as a map. Where it is large --- contact
and alignment, or the first frame of an episode where the intended goal is hidden --- a longer
commitment is priced higher; where it is near zero the observation suffices and short commitments tie.
The map is the critic-side counterpart of the policy-side non-Markovianity metric of prior
work~\citep{lazzati2026chunking}; verifying that it tracks genuine occlusion and aleatoric structure,
rather than an arbitrary temporal correlation, is the check that the history conditioning of
\eqref{eq:rcv-input} carries state and not a shortcut.
